\subsection{The factor \texorpdfstring{$12$}{12} and the Bott--Chern first Chern class}

Let
\[
 \exp\colon H^1(X,\mathcal O_X)\longrightarrow\operatorname{Pic}(X)
\]
be induced by $z\mapsto e^{2\pi iz}$.  By
\cref{prop:leray-e2},
\begin{equation}
 H^0(B,R^1f_*\mathbb Z)=\mathbb Z(12\gamma).
\label{eq:bcpr-global-integral-line}
\end{equation}
The exponential--Leray comparison gives a line bundle $M$ on $B$ such
that
\begin{equation}
 \exp(12\mathfrak g)=f^*M,
 \qquad c_1(M)=d_2(12\gamma).
\label{eq:bcpr-exponential-transgression}
\end{equation}
For completeness, the first assertion follows because
\cref{eq:bcpr-global-integral-line} says that the restriction of
$12\mathfrak g$ to every fibre lies in the image of integral cohomology
and hence exponentiates to the trivial bundle.  The low-degree Leray
sequence for $\mathcal O^*$ then makes the resulting bundle a pullback.
On a disc--pair-of-pants cover of $B$, choose local integral lifts of
$12\gamma$.  Their differences on the annular overlaps form a Cech
one-cochain, and its Cech coboundary is simultaneously the two-cocycle
defining the Leray transgression and the two-cocycle of logarithms defining
$c_1(M)$.  This proves the second assertion with the same sign convention;
it is the exponential-sequence version of the local-lift calculation in
the proof of \cref{lem:first-transgression}.

Orient the generator $\omega\in H^2(B;\mathbb Z)$ by
\begin{equation}
 \omega=c_1(\mathcal O_B(1)).
\label{eq:bcpr-base-orientation}
\end{equation}
Then \cref{lem:first-transgression} and
\cref{eq:bcpr-exponential-transgression} give
\begin{equation}
 \exp(12\mathfrak g)=f^*\mathcal O_B(1).
\label{eq:bcpr-factor-twelve}
\end{equation}
If the opposite orientation is used for the primitive generator in
\cref{lem:first-transgression}, the corresponding statement is
$\exp(-12\mathfrak g)=f^*\mathcal O_B(1)$.  Thus the canonical integral
data fixes the absolute factor $12$, while its sign is the orientation
choice fixed in \cref{eq:bcpr-base-orientation}.

Put
\begin{equation}
 a=\partial G,
 \qquad a^\dagger=\bar\partial\bar G,
 \qquad
 \mathsf A=[a]_{BC},
 \qquad \mathsf B=[a^\dagger]_{BC}.
\label{eq:bcpr-a-b-definitions}
\end{equation}
Both forms are $d$-closed.  We now compute their Chern-class normalization
without choosing local metric weights.  For $z\in\mathbb C$, equip the
smooth trivial line with a global frame $s$ and the holomorphic structure
\begin{equation}
 \bar\partial_zs=-2\pi i zG\otimes s.
\label{eq:bcpr-holomorphic-structure}
\end{equation}
It is integrable because $\bar\partial G=0$.  If
$G=\bar\partial\phi_i$ on a good cover, then
$e_i=e^{2\pi iz\phi_i}s$ is $\bar\partial_z$-holomorphic and
\[
 e_j=e^{2\pi iz(\phi_j-\phi_i)}e_i.
\]
The holomorphic cocycle $\phi_j-\phi_i$ represents $\mathfrak g$ under
the Cech--Dolbeault comparison, so
\cref{eq:bcpr-holomorphic-structure} is precisely the line bundle
$\exp(z\mathfrak g)$ for the displayed exponential convention.

Give $s$ unit length.  The unit-metric Chern connection is
\begin{equation}
 \nabla_z=d-2\pi i\bigl(zG+\bar z\,\bar G\bigr).
\label{eq:bcpr-unit-chern-connection}
\end{equation}
Indeed, its $(0,1)$ part is
\cref{eq:bcpr-holomorphic-structure}, and its connection one-form is
purely imaginary, so it is metric compatible; uniqueness identifies it
with the Chern connection.  Since $\bar\partial G=0$ and
$\partial\bar G=0$, its curvature is the global $(1,1)$ form
\begin{equation}
 F_{\nabla_z}
 =-2\pi i\bigl(z\,\partial G+
                 \bar z\,\bar\partial\bar G\bigr)
 =-2\pi i\bigl(za+\bar z\,a^\dagger\bigr).
\label{eq:bcpr-unit-chern-curvature}
\end{equation}
With the Bott--Chern convention
\begin{equation}
 c_{1,BC}(L,h)=
 \left[\frac{i}{2\pi}F_{\nabla^h}\right]_{BC}
\label{eq:bcpr-chern-convention}
\end{equation}
of \cite[\S8.a]{Schweitzer2007},
\cref{eq:bcpr-unit-chern-curvature} gives, for every $z\in\mathbb C$,
\begin{equation}
 c_{1,BC}\bigl(\exp(z\mathfrak g)\bigr)
 =z\mathsf A+\bar z\,\mathsf B.
\label{eq:bcpr-chern-exp-identity}
\end{equation}
Consequently, if $\varepsilon=1$ for the orientation fixed by
\cref{eq:bcpr-base-orientation} and $\varepsilon=-1$ for its opposite,
then
\begin{equation}
 c_{1,BC}\bigl(f^*\mathcal O_B(1)\bigr)
 =12\varepsilon\bigl(\mathsf A+\mathsf B\bigr).
\label{eq:bcpr-c1-pullback}
\end{equation}
This is the convention-consistent form of the factor-$12$ identity; no
analytic rescaling of the Picard coordinate changes its integral content.

